\section{Assumptions}\label{assum}
In the design of \texttt{Blossom}, we maintain certain assumptions regarding the behavior of computer systems, as they relate to the utility of our protocol. These assumptions are both general and specific, relating to both the hardware and software required to implement our system. Operating with these assumptions we are able to ignore some quotidian aspects of our technology, as well as implement technologies that have been previously untested. We have utilized assumptions during both the development and analysis of our protocol, by means which a reasonable and accounted for at each relevant occasion. To clarify which assumptions the authors employed in the drafting of this document, we will outline our assumptions in the following points:
\begin{description}
\item[Byzantine Behavior] As previously defined as equivalent to ‘arbitrary’ behavior, can be expected to occur at some generalized rate in the natural world. While we have tested our protocol in what is equivalent to laboratory conditions, we suppose the typical fault occurrence of one fault for every \(10^3\) nodes every two weeks.
\item[Protocol Security] Security is of utmost importance, and while this paper may not describe the methods developers may implement to mitigate the risks related to \texttt{Blossom}, we recognize that such solutions do exist (e.g. ZK-proofs, financial incentives, and or supplementary consensus protocols) which may resolve said risks.
\item[Fault Disincentives] Another topic that we do not directly reference in this document, which would also be valuable in the mitigation of possibly byzantine behavior. Such disincentives would limit or remediate the perpetual misbehavior of select nodes while providing other network members the ability to act on said issues. This and other ‘social issues’ are not pertinent to the design of this protocol, but pertinent to the implementation of this protocol, thus we believe some such system must be implemented along with the public use of \texttt{Blossom}.
\item[Share Credentials] An important characteristic of decentralized systems, implementations of which have been proven as both effective and efficient across scales. Recognizing the viability of such an architecture, for the purpose of our development timeline, we chose to utilize a centralized database for actions requiring peer lookup. While we do not expect this behavior to continue beyond this initial proof-of-concept, we expect the use of a database should be negligible to the performance of our protocol.
\item[State Determinism] Best defined as a belief that two similar machines running the same software can independently reach an identical state, given both machines are provided the same input data. This assumption requires the data provided to said nodes to be identical in every way (i.e. indistinguishable).
\end{description}
